\documentclass[%
 reprint,
%superscriptaddress,
%groupedaddress,
%unsortedaddress,
%runinaddress,
%frontmatterverbose, 
%preprint,
%preprintnumbers,
%nofootinbib,
%nobibnotes,
%bibnotes,
 amsmath,amssymb,
 aps,
%pra,
%prb,
%rmp,
%prstab,
%prstper,
%floatfix,
]{revtex4-2}

\usepackage{graphicx}% Include figure files
\usepackage{dcolumn}% Align table columns on decimal point
\usepackage{bm}% bold math
\usepackage{hyperref}% add hypertext capabilities
\usepackage{amsthm}
\newcommand{\E}{\mathcal{E}}
\newcommand{\norm}[1]{\left\lVert#1\right\rVert}
%\usepackage[mathlines]{lineno}% Enable numbering of text and display math
%\linenumbers\relax % Commence numbering lines

%\usepackage[showframe,%Uncomment any one of the following lines to test 
%%scale=0.7, marginratio={1:1, 2:3}, ignoreall,% default settings
%%text={7in,10in},centering,
%%margin=1.5in,
%%total={6.5in,8.75in}, top=1.2in, left=0.9in, includefoot,
%%height=10in,a5paper,hmargin={3cm,0.8in},
%]{geometry}

\newtheorem{theorem}{Theorem}[section]
\newtheorem{corollary}{Corollary}[theorem]
\newtheorem{lemma}[theorem]{Lemma}
\newtheorem{definition}{Definition}[section]

\begin{document}

\preprint{APS}

\title{Learning a Quantum Channel}% Force line breaks with \\

\author{Xiang Zou}
\email{xiang.zou@mail.utoronto.ca}
\affiliation{%
 Department of Physics, University of Toronto, 60 St. George Street, Toronto ON, M5S 1A7, Canada
}%
 %\altaffiliation[Also at ]{Physics Department, XYZ University.}%Lines break automatically or can be forced with \\
\author{Yuxuan Zhang}%
%\email{hklo@ece.utoronto.ca}
\affiliation{%
 Department of Physics, University of Toronto, 60 St. George Street, Toronto ON, M5S 1A7, Canada
}%
\author{collaborators}%
\affiliation{%
 Department of Physics, University of Toronto, 60 St. George Street, Toronto ON, M5S 1A7, Canada
}%
\date{\today}% It is always \today, today,
             %  but any date may be explicitly specified

\begin{abstract}
In this paper, we would like to discuss the theory and simulation results behind learning a quantum Channel.
\end{abstract}

%\keywords{Suggested keywords}%Use showkeys class option if keyword
                              %display desired
\maketitle

%\tableofcontents

\section{Introduction}

Quantum Channel Learning is of great theoretical interest, aiming to learn an unknown quantum channel described by a CPTP (Completely Positive Trace Preserving) map.

\section{Theoretical Analysis}
Several results discuss the bounds of using Quantum Machine Learning Models (QMLM) to study a general unitary operation on several input modes. We would like to generalize the idea of learning a general unitary operation to learning a general (unital) quantum channel. That we have to take into consideration dephasing and decoherence effects.

\subsection{Problem Statement}
Following previous results in learning a random unitary operation, we consider a quantum machine learning model (QMLM) as a parameterized quantum channel, i.e., a completely positive trace-preserving (CPTP) map that is parameterized. This QMLM is denoted as $\mathcal{E}^{QMLM}_{\alpha}(\cdot)$, where $\alpha$ denotes the set of parameters that can be tuned in our parameterized QMLM channel.

Then, given a target quantum Channel $\mathcal{E}$, and a set of known input and output $\{\rho_i\}$ and $\{\mathcal{E}(\rho_i)\}$, we would like to construct a QMLM model $\E^{QMLM}_{\alpha}$, such that the diamond distance between the 2 channel
\begin{equation}
    \norm{\E-\E^{QMLM}_{\alpha}}_\diamond
\end{equation}
is sufficiently small, or is bounded by a constant. Here, just a reminder that the diamond distance between 2 channels is defined as
\begin{equation}
    \norm{\E_1 - \E_2}_{\diamond} = \sup_{\rho} \left\| (\E_1 \otimes \operatorname{Id}_k)(\rho) - (\E_2 \otimes \operatorname{Id}_k)(\rho) \right\|_1
\end{equation}

To train a QMLM capable of this task, we need a well-defined and easily computable cost function for the machine-learning implementation. Unlike the task of learning a general unitary operation, learning a general quantum channel is much harder, as defining the cost functions is not as easy. The easiest approach is to use the diamond distance between the channels to define the cost function. This approach is very hard to implement as the 1-norm is usually very hard to compute, not to mention finding the diamond distance between 2 channels, which requires finding the supremum of the 1-norms. In order to get a scheme that is easy to implement, here we have summarized three different approaches to learn a general unital quantum channel.

\subsection{Method 1: Learning the Purification}

The easiest and most trivial approach might be that, since we know that every quantum channel described by a completely positive trace-preserving (CPTP) map can be realized as a unitary evolution on a larger Hilbert space, involving the original system and an auxiliary (ancilla) system. Also, the purification is not unique for a given channel.

Hence, if we can learn one of the purification methods of the quantum channel, which is a unitary operation that is well-constructed by previous studies, then we can learn the $\E^{QMLM}_\alpha$ channel of the target channel. 

This is the simplest approach, while it has obvious drawbacks. In previous studies, to learn a general unitary, the cost function is usually defined as 
\begin{equation}
\left.C_{\mathcal{D}_{\mathcal{Q}}(N)}(\boldsymbol{\alpha})=1-\frac{1}{N} \sum_{j=1}^N\left|\left\langle\Psi^{(j)}\right| V(\boldsymbol{\alpha})^{\dagger} U\right| \Psi^{(j)}\right\rangle\left.\right|^ 2
\end{equation}

Where calculating $V(\boldsymbol{\alpha})^{\dagger} U$ is essential. And as we don't have access to the ancillary information of the purification, it would be almost impossible to calculate $V(\boldsymbol{\alpha})^{\dagger} U$, and thus this method only works if we guess the purification at the start. 

\subsection{Method 2: Using the Channel Fidelity}

The problem is that calculating the Channel Fidelity is very hard in general. While calculating the 2-norm is, in general, way easier.

The fidelity of 2 quantum states is defined as 
\begin{equation}
    F(\rho, \sigma) = \left( \operatorname{Tr} \sqrt{ \sqrt{\rho} \sigma \sqrt{\rho} } \right)^2
\end{equation}

The main idea of this approach is to use $F(\E(\rho), \E^{QMLM}_\alpha(\rho))$ as the cost function, and formulate the QMLM that minimizes the cost function. The main challenge is that the fidelity is, in general, hard to compute, and we need to figure out the distribution of the training dataset.

\subsection{Method 3: Using Pauli Transfer Matrix}

We denote the Choi-representation of the quantum channel $\E$ as $J(\E)$. And we have the following inequality from the norm inequality.

\begin{equation}
    \norm{J(\E_1)-J(\E_2)}_2\le \norm{\E_1-\E_2}_\diamond \le \sqrt{D}\norm{J(\E_1)-J(\E_2)}_2
\end{equation}
Where $D$ is the Matrix dimension of the channel. The upper bound is not tight, for practical reasons and average case calculation, bounding $\norm{J(\E_1)-J(\E_2)}_2$ should be sufficient.

The Pauli Transfer Matrix (PTM) $R$ is the matrix representation of a quantum channel $\mathcal{E}$ in the basis of Pauli operators. For an $n$-qubit quantum channel, the PTM is a $4^n \times 4^n$ real-valued matrix whose elements $R_{ij}$ are defined as:

\begin{equation}
    R_{ij} = \frac{1}{d} \mathrm{Tr}\left[ P_i^\dagger \mathcal{E}(P_j) \right]
\end{equation}

Since the Pauli operators are Hermitian ($P_i^\dagger = P_i$), this is equivalent to:

\begin{equation}
    R_{ij} = \frac{1}{d} \mathrm{Tr}\left[ P_i \mathcal{E}(P_j) \right]
\end{equation}

where:
\begin{description}
    \item[$R$] is the Pauli Transfer Matrix.
    \item[$\mathcal{E}$] is the quantum channel, which is a completely positive and trace-preserving (CPTP) linear map.
    \item[$P_i, P_j$] are operators from the $n$-qubit Pauli basis. This basis consists of all $4^n$ tensor products of the single-qubit Pauli matrices $\{I, \sigma_x, \sigma_y, \sigma_z\}$.
    \item[$d$] is the dimension of the Hilbert space, where $d=2^n$ for a system of $n$ qubits.
    \item[${\mathrm{Tr}[\cdot]}$] denotes the trace of a matrix.
\end{description}

Since both the Choi-representation and the PTM are equivalent ways to represent quantum channels. Hence, we have that the Frobenius Norms,
\begin{equation}
    \norm{J(\E_1)-J(\E_2)}_F^2=\norm{R(\E_1)-R(\E_2)}_F^2
\end{equation}

Therefore, we can define our metric as below, where we denote the function that we are calculating as $l(\alpha, a, b) = \left(\mathrm{Tr}\left[P_a(\E-\E^{QMLM}_\alpha)P_b\right]\right)^2$
\begin{equation}
\begin{split}
    F(\E, \E^{QMLM}_\alpha) &= \norm{R(\E)-R(\E^{QMLM}_\alpha)}^2\\
    &=\frac{1}{d}\sum_{a=0}^{d-1}\sum_{b=0}^{d-1}\left(\mathrm{Tr}\left[P_a(\E-\E^{QMLM}_\alpha)P_b\right]\right)^2\\
    &=\mathbb{E}\left[\left(\mathrm{Tr}\left[P_a(\E-\E^{QMLM}_\alpha)P_b\right]\right)^2\right]\\
    &=E [Tr(\E P_{in}-\E^{QMLM} P_{in})^2]
\end{split}
\end{equation}

There are exponentially many different Pauli strings $P_a, P_b$; hence, calculating the exact $F(\E, \E^{QMLM}_\alpha)$ is exponentially expensive. Thus, we would like to approximate the value of $F(\E, \E^{QMLM}_\alpha)$ using some finite set of samples, which is shown below. Suppose we have a subset of all possible Pauli strings $P_N\subset\{P_i\}_{i=0}^{d}$, and $|P_N|=N$, then,
\begin{equation}
\begin{split}
    F'_{P_N}(\E, \E^{QMLM}_\alpha)
    =\frac{1}{N}\sum_{a=0}^{N}\sum_{b=0}^{N}\left(\mathrm{Tr}\left[P_a(\E-\E^{QMLM}_\alpha)P_b\right]\right)^2
\end{split}
\end{equation}

Then, we can define the generalization bounds as below, determined by how well the $F'_{P_N}$ with a small sample size $N$ reflects the overall distribution $F$.

\begin{equation}
    gen_{P_N}=F(\E, \E^{QMLM}_\alpha) - F'_{P_N}(\E, \E^{QMLM}_\alpha)
\end{equation}

\section{Prediction error bounds for quantum machine learning models}
In this section, we would like to calculate the bounds for the generalization error $gen_{P_N}$ for a given set of randomly chosen Pauli Strings $P_N$ and $T$ independent parameterized 2-qubit $\mathcal{CPTP}$ maps.

We first quote 2 lemmas and a theorem proven in previous papers.

\begin{lemma}
    (Covering number bounds for 2-qubit unitaries). Let $\norm{\cdot}$ be a unitarily invariant norm on complex $4 \times 4$ matrices. The covering number of the set of 2-qubit unitaries $U(\mathbf{C}^2 \otimes \mathbf{C}^2)$ w.r.t. the norm $\norm{\cdot}$ can be bounded as
    \begin{equation}
    \mathcal{N}\left(\mathcal{U}\left(\mathbb{C}^2 \otimes \mathbb{C}^2\right),\|\cdot\|, \varepsilon\right) \leqslant\left(\frac{6\left\|I_{\mathbb{C}^4}\right\|}{\varepsilon}\right)^{32}, \quad \text { for } 0<\varepsilon \leqslant\left\|I_{\mathbb{C}^4}\right\|
    \end{equation}
\end{lemma}

\begin{lemma}
    (Covering number bounds for 2-qubit CPTP maps). The covering number of the set of 2-qubit CPTP maps $\mathcal{CPTP}(\mathbf{C}^2 \otimes \mathbf{C}^2)$ w.r.t. the diamond distance can be bounded as
    \begin{equation}
    \mathcal{N}\left(\mathcal{CPTP}\left(\mathbb{C}^2 \otimes \mathbb{C}^2\right),\|\cdot\|_\diamond, \varepsilon\right) \leqslant\left(\frac{6}{\varepsilon}\right)^{512}, \quad \text { for } 0<\varepsilon \leqslant1
    \end{equation}
\end{lemma}

\begin{theorem}
    (Metric entropy bounds for QMLMs of CPTP maps). Let $\E^{QMLM}_\theta(\cdot)$ be an n-qubit QMLM with T parameterized 2-qubit CPTP maps and an arbitrary number of non-trainable, global CPTP maps. Let $\mathcal{CPTP}^{QMLM}\subset \mathcal{CPTP}((C^2)^{\otimes n})$ denote the set of n-qubit CPTP maps that can be implemented by the circuit QMLM $\E^{QMLM}_\theta(\cdot)$.

    For any $\varepsilon \in (0, 1]$, there exists an $\varepsilon$-covering net $\mathcal{N}_\varepsilon$ of $\mathcal{CPTP}^{QMLM}$ w.r.t. the diamond distance such that the logarithm of its size can be upper bounded as
    \begin{equation}
        \log(|\mathcal{N_\varepsilon}|)\le 512T\log\left(\frac{6T}{\varepsilon}\right)
    \end{equation}
    \label{thm: Metric entropy bounds}
\end{theorem}

Now we would like to calculate the bounds. First, we will begin with a technically simple proof of a generalization bound for a fixed-architecture QMLM, and subsequently, we will refine this bound.

\begin{theorem}
    (Prediction error bound for quantum machine learning - Fixed structure (Preliminary version)). Let $\E^{QMLM}_\theta(\cdot)$ be a QMLM with a fixed architecture consisting of $T$ parameterized 2-qubit CPTP maps and an arbitrary number of non-trainable, global CPTP maps. Let $P$ be a probability distribution over input-output pairs. Suppose that, given training data $S = \{(P_i, P_j)\}_{i=1}^N$ of size N, our optimization yields the parameter setting $\theta^* = \theta^*(S)$. Then, with probability at least $ 1-\delta$ over the choice of i.i.d. training data $S$ of size $N$ according to $P$,
    \begin{equation}
        F\left(\boldsymbol{\theta}^*\right)-F'_S\left(\boldsymbol{\theta}^*\right) \in \mathcal{O}\left(C_{\operatorname{loss}}\left(\sqrt{\frac{T \log (T N)}{N}}+\sqrt{\frac{\log (1 / \delta)}{N}}\right)\right)
    \end{equation}
\end{theorem}

\begin{proof}
    Here, the independent random variable that we are considering is $l(theta, P_a, P_b)=\left(\mathrm{Tr}\left[P_a(\E-\E^{QMLM}_\theta)P_b\right]\right)^2$, this random variable lies between $[-c, c]$. Hence, according to Hoeffding's contraction inequality, 
    \begin{equation}
        \mathbb{P}_S\left[F'_S(\theta)-F(\theta)> \eta \right] \le exp\left(-\frac{N\eta^2}{2c^2}\right)
    \end{equation}

    Next, we let $\varepsilon=\sqrt{T/N} >0$ and let $\mathcal{N}_\varepsilon$ be the $\varepsilon$-covering net of the set of CPTP maps that can be implemented by the QML.  We take the union bound over $\mathcal{N}_\varepsilon$, that we obtain:
    \begin{equation}
        \mathbb{P}_S\left[\exists \E^{QMLM}_\theta:F'_S(\theta)-F(\theta)> \eta \right] \le |\mathcal{N}_\varepsilon| \cdot exp\left(-\frac{N\eta^2}{2c^2}\right)
    \end{equation}

    As we took $N_\varepsilon$ to be an $\varepsilon$-covering net (w.r.t. the diamond norm) of the class of CPTP maps that the QMLM can implement, and since $\norm{\E-\E^{CPTP}_\theta}_\diamond$ directly implies, for all pauli strings $P_i, P_j$,
    \begin{equation}
        |F'_S-F|\le\norm{\E-\E^{CPTP}_\theta}_\diamond\le c\varepsilon
    \end{equation}
    Hence, the optimal parameter setting $\theta^*$ will relate as follows,
    \begin{equation}
    \begin{split}
        \mathbb{P}_S\left[F'_S(\theta^*)-F(\theta^*)> \eta +2\varepsilon \right] &\le \mathbb{P}_S\left[\exists \E^{QMLM}_\theta:F'_S(\theta)-F(\theta)> \eta \right]\\
        &\le |\mathcal{N}_\varepsilon| \cdot exp\left(-\frac{N\eta^2}{2c^2}\right)
    \end{split}
    \end{equation}

    Therefore, for any $\delta \in (0, 1)$, and we set $\eta = c\sqrt{\frac{\log{(|\mathcal{N}_\varepsilon|/\delta})}{N}}$. Then, after simple calculation, we can get that, with the probability of $1-\delta$ over the choice of training data $S$ of size $N$, we have 
    \begin{equation}
        F'_S(\theta^*)-F(\theta^*)\le c\sqrt{\frac{\log{(|\mathcal{N}_\varepsilon|/\delta})}{N}}+2c\sqrt{\frac{T}{N}}
    \end{equation}

    According to the theorem~\ref{thm: Metric entropy bounds}, we have $\log{|\mathcal{N}_\varepsilon|}\le 512 T\log{\left(
    \frac{6T}{\varepsilon}\right)
    }$
    Hence, a simple calculation leads to
\begin{equation}
        F\left(\boldsymbol{\theta}^*\right)-F'_S\left(\boldsymbol{\theta}^*\right) \in \mathcal{O}\left(C_{\operatorname{loss}}\left(\sqrt{\frac{T \log (T N)}{N}}+\sqrt{\frac{\log (1 / \delta)}{N}}\right)\right)
    \end{equation}
\end{proof}

\section{Training with weight-1 Pauli String}

In our quantum channel learning model, we do the training primarily using the Pauli String of weight-1, and we believe it will be a good learning set.

\subsection{When does a channel preserve operator products?}

This subsection clarifies when one may replace $\mathcal{C}(AB)$ by $\mathcal{C}(A)\mathcal{C}(B)$ at the
operator level.

A (finite-dimensional) quantum channel is a completely positive,
trace-preserving (CPTP) linear map
\(\mathcal{C}:\mathcal{B}(\mathcal{H})\to\mathcal{B}(\mathcal{H})\),
where $\mathcal{B}(\mathcal{H})$ denotes the bounded operators on the system Hilbert space $\mathcal{H}$. A unitary channel is given by
\begin{equation}
  \mathcal{U}(X) = U X U^{\dagger},
  \qquad X\in\mathcal{B}(\mathcal{H}),
\end{equation}
for a unitary operator $U$ on $\mathcal{H}$. The map $\mathcal{U}$ is linear, completely positive, trace-preserving
and unital, i.e. $\mathcal{U}(\openone)=\openone$.

\begin{theorem}
Let $\mathcal{U}(X)=UXU^{\dagger}$ be a unitary channel. Then
$\mathcal{U}$ is a unital $*$-automorphism of the operator algebra
$\mathcal{B}(\mathcal{H})$: for all $A,B\in\mathcal{B}(\mathcal{H})$,
\begin{equation}
  \mathcal{U}(AB)=\mathcal{U}(A)\,\mathcal{U}(B),
  \qquad
  \mathcal{U}(A^{\dagger})=\mathcal{U}(A)^{\dagger}.
\end{equation}
In particular, if $A$ and $B$ are Pauli strings then
$\mathcal{U}(AB)=\mathcal{U}(A)\mathcal{U}(B)$.
\end{theorem}

\begin{proof}
Linearity is immediate from the definition. For the adjoint we have
\begin{equation}
  \mathcal{U}(A^{\dagger})
  = U A^{\dagger} U^{\dagger}
  = (UAU^{\dagger})^{\dagger}
  = \mathcal{U}(A)^{\dagger}.
\end{equation}
For multiplicativity, we compute
\begin{equation}
  \mathcal{U}(AB)
  = U (AB) U^{\dagger}
  = (UAU^{\dagger})(UBU^{\dagger})
  = \mathcal{U}(A)\,\mathcal{U}(B),
\end{equation}
where we used $U^{\dagger}U=\openone$ to combine the middle factors.
Thus $\mathcal{U}$ is a $*$-homomorphism of $\mathcal{B}(\mathcal{H})$
onto itself, i.e.\ a $*$-automorphism.
\end{proof}

Therefore, whenever the channel describing the evolution is implemented by conjugation with a single unitary, products of operators (and in particular products of Pauli strings) are preserved under the channel in the sense that $\mathcal{U}(AB)=\mathcal{U}(A)\mathcal{U}(B)$.

Thus, limiting to only learning a unital channel, all different Pauli strings can be written as the product of one or several weight-1 Pauli strings. Therefore, learning how the weight-1 Pauli string changes under the unital channel is sufficient to capture the effect of the channel on all Pauli strings.

However, a general CPTP map $\mathcal{E}$ admits a Kraus representation
\begin{equation}
  \mathcal{E}(X) = \sum_{k} K_k X K_k^{\dagger},
  \qquad \sum_{k} K_k^{\dagger}K_k = \openone.
\end{equation}
Unless there is only a single Kraus operator and it is unitary, such a map is not a $*$-homomorphism, and in general
\begin{equation}
  \mathcal{E}(AB) \neq \mathcal{E}(A)\,\mathcal{E}(B)
\end{equation}
even when $A$ and $B$ are Pauli strings.
This can be seen explicitly from
\begin{equation}
  \mathcal{E}(AB)
  = \sum_{k} K_k A B K_k^{\dagger},
  \quad
  \mathcal{E}(A)\mathcal{E}(B)
  = \sum_{k,\ell} K_k A K_k^{\dagger} K_{\ell} B K_{\ell}^{\dagger},
\end{equation}
whose cross terms with $k\neq \ell$ do not cancel in general.

As a concrete example, consider the single-qubit dephasing channel
\begin{equation}
  \mathcal{D}_p(\rho) = p\,\rho + (1-p)\,Z\rho Z,
  \qquad 0<p<1,
\end{equation}
which is CPTP but not unitary for $0<p<1$.
Taking $A=B=X$, one finds
\begin{equation}
  \mathcal{D}_p(XX) = \mathcal{D}_p(\openone) = \openone,
  \qquad
  \mathcal{D}_p(X) = (2p-1)X,
\end{equation}
so that
\begin{equation}
  \mathcal{D}_p(X)\,\mathcal{D}_p(X) = (2p-1)^2 \openone \neq \openone
\end{equation}
whenever $0<p<1$.
Thus, for this non-unitary channel,
$\mathcal{D}_p(AB)=\mathcal{D}_p(A)\mathcal{D}_p(B)$ already fails for
the simple case $A=B=X$.

\paragraph{Multiplicative domain viewpoint.}
Given a completely positive map $\mathcal{E}$, its \emph{multiplicative
domain} is the set
\begin{equation}
\begin{split}
    \mathrm{MD}(\mathcal{E})
  := \Bigl\{A \in \mathcal{B}(\mathcal{H}) \,\Big|&\,
  \mathcal{E}(A^{\dagger}A)
  = \mathcal{E}(A^{\dagger})\mathcal{E}(A)
  \ \text{ and }\\
  &\mathcal{E}(AA^{\dagger})
  = \mathcal{E}(A)\mathcal{E}(A^{\dagger})\Bigr\}
\end{split}
\end{equation}
It is a $C^{*}$-subalgebra of $\mathcal{B}(\mathcal{H})$, and
$\mathcal{E}$ is a $*$-homomorphism when restricted to
$\mathrm{MD}(\mathcal{E})$.
In particular, for any $A\in\mathrm{MD}(\mathcal{E})$ and all $B$,
\begin{equation}
  \mathcal{E}(AB) = \mathcal{E}(A)\,\mathcal{E}(B),
  \qquad
  \mathcal{E}(BA) = \mathcal{E}(B)\,\mathcal{E}(A).
\end{equation}
For a unitary channel $\mathcal{U}(X)=UXU^{\dagger}$ one has
$\mathrm{MD}(\mathcal{U})=\mathcal{B}(\mathcal{H})$, so multiplicativity holds for all operators (and hence for all Pauli strings). In contrast, for a generic noisy channel, the multiplicative domain is a proper subalgebra, and Pauli strings do not, in general, lie in
$\mathrm{MD}(\mathcal{E})$.

In summary, the identity $\mathcal{C}(AB)=\mathcal{C}(A)\mathcal{C}(B)$ for all Pauli strings $A,B$ is valid precisely when $\mathcal{C}$ is a unitary channel (or, more generally, a $*$-homomorphism), and it fails for generic non-unitary CPTP maps.




\section{Numerical Result}
Here we would like to show the numerical result.

Can we consider a QKD channel? Probably, people will be interested. Thought about how to excite physicists.


\section{Possible setups}
\begin{enumerate}
    \item learning noisy random circuits (should be easier to learn when considering deep circuits - constraint on a given (all-0) input state?)
    \item compressing open quantum dynamics, such as Lindbladian dynamics
    \item matching local observables (can use a classical shadow trick in this case)
\end{enumerate}


\bibliography{apssamp}% Produces the bibliography via BibTeX.

\end{document}
%
% ****** End of file apssamp.tex ******
